\documentclass{subfiles}
\begin{document}
    \begin{definition*}
        Let \(\C* \subseteq A^{n}\) be a code. If we have that \(\C*\) is 
            a vector subspace of \(A^{n}\), we say that \(\C*\) is linear.
            That is, \(\C*\) is a linear code if:
            \begin{enumerate}
                \item For any \(\bv{c}_{1}, \bv{c}_{2} \in \C*,
                    \bv{c}_{1} + \bv{c}_{2} \in C\); and

                \item For any \(\bv{c} \in \C*, \forall \alpha \in A, 
                    \alpha\bv{c} \in \C*\).
            \end{enumerate}
    \end{definition*}

    Each code \(\C*\) is identified by the tuple \((n, k, d)\), 
        where \(n\) is the size of the words (i.e. the number of components),
        \(k\) the dimesions of \(\C*\) and \(d\) the minimal distance.
        In the case of linear codes we write \([n, k, d]\) instead of \((n, k, d)\).

    The question now is: how do we can represent a linear code?
        Since \(\C*\) is a vector subspace, it must have a basis of vectors that 
        generates the whole subspace. If \(\bv{c}_{1}, \ldots, \bv{c}_{k}\)
        form a basis for \(\C*\), we can denote by 
        \[
            G = \begin{pmatrix}
                \bv{c}_{1} \\ 
                \vdots     \\ 
                \bv{c}_{k} \\
            \end{pmatrix}
        \]
        the \emph{generator matrix} of \(\C*\). It can be easely seen that any 
        word in \(\C*\) is given by the linear combination of rows in \(G\).

    Alternatively, we can make use of a \emph{parity-check matrix}.
        That is an \(r \times n\) matrix \(H\) over \(A\), such that 
        \[
            \bv{c} \in \C* \iff H\bv{c}^{T} = \bv{0}.
        \]
        Note that \(r = \rank H = n - k\), from the dimesional relationship.

    \subsubsection{Theorems and inequalities}
    \subfile{../subsub/6.2.1 - Theorems and inequalities}

    \subsubsection{Hmming codes}
    \subfile{../subsub/6.2.2 - Hamming codes}

    \subsubsection{Hadamard codes}
    \subfile{../subsub/6.2.3 - Hadamard codes}

\end{document}
